‫\definecolor{dsdlcodebg}{RGB}{232,238,245}
‫\definecolor{dsdlcodeborder}{RGB}{208,215,222}
‫\newsavebox{\dsdlcodebox}
‫\newenvironment{dsdlcodeblock}
‫{%
‫  \par\noindent
‫  \begin{latin}
‫  \begin{lrbox}{\dsdlcodebox}
‫  \begin{minipage}{1.1\linewidth}
‫  \small
‫}
‫{%
‫  \end{minipage}
‫  \end{lrbox}
‫  \setlength{\fboxsep}{6pt}
‫  \fcolorbox{dsdlcodeborder}{dsdlcodebg}{\usebox{\dsdlcodebox}}
‫  \end{latin}
‫  \par
‫}
‫
‫% -------------------------------------------------------
‫%  Abstract
‫% -------------------------------------------------------
‫
‫\section{مقدمه}
‫
‫ضرب مستقیم دودویی به روش «شیفت و جمع» در هر مرحله یک بیت ضریب را بررسی می‌کند و برای
‫یک ضریب $N$ بیتی به $N$ مرحله‌ی محاسباتی نیاز دارد. روش بوث با بررسی دو بیت متوالی ضریب،
‫دنباله‌های یک را به تفاضل دو توان از دو تبدیل می‌کند و ضرب اعداد علامت‌دار مکمل دو را با یک
‫فرایند یکنواخت انجام می‌دهد. این روش نخستین بار توسط اندرو بوث معرفی شد.
‫
‫هدف آزمایش پنجم، طراحی یک ضرب‌کننده‌ی بوث با جداسازی مسیر داده\پاورقی{Datapath} و واحد
‫واپایش\پاورقی{Control Unit} است. صورت آزمایش همچنین تصریح می‌کند که شیفت‌دهنده باید بتواند
‫در یک نبض ساعت بیش از یک بیت به راست شیفت دهد. طرح حاضر این شرط را با یک شیفت‌دهنده‌ی
‫بشکه‌ای\پاورقی{Barrel Shifter} ترکیبی و یک نقشه‌ی مرز بوث برآورده می‌کند.
‫
‫
‫
‫
‫\section{مبانی الگوریتم بوث}
‫
‫در الگوریتم بوث، ثبات $Q$ ضریب و بیت $Q_{-1}$ مقدار قبلی کم‌ارزش‌ترین بیت ضریب را نگه
‫می‌دارد. در آغاز $Q_{-1}=0$ است. جدول~\رجوع{جدول:عملیات بوث} عملیات متناظر با زوج
‫${Q_0,Q_{-1}}$ را نشان می‌دهد.
‫
‫\شروع{لوح}[ht]
‫\تنظیم‌ازوسط
‫\شرح{انتخاب عملیات در الگوریتم بوث}
‫\شروع{جدول}{|c|c|}
‫\خط‌پر
‫\سیاه زوج ${Q_0,Q_{-1}}$ & \سیاه عملیات روی انباره‌ی $A$ \\
‫\خط‌پر \خط‌پر
‫$00$ & بدون تغییر \\
‫$01$ & $A \leftarrow A+M$ \\
‫$10$ & $A \leftarrow A-M$ \\
‫$11$ & بدون تغییر \\
‫\خط‌پر
‫\پایان{جدول}
‫\برچسب{جدول:عملیات بوث}
‫\پایان{لوح}
‫
‫پس از عملیات جدول، ثبات ترکیبی ${A,Q,Q_{-1}}$ به صورت حسابی به راست شیفت می‌خورد.
‫برای نمونه، $7=00000111_2$ را می‌توان به صورت $8-1$ نوشت؛ بنابراین ضرب $M$ در ۷ به
‫جای سه جمع متوالی، با یک تفریق در ابتدای دنباله‌ی یک‌ها و یک جمع در انتهای آن نمایش داده
‫می‌شود. این تفسیر مبنای پرش چندبیتی در طرح حاضر است.
‫
‫
‫\section{معماری کل سامانه}
‫
‫پیکربندی فعلی با پارامتر $N=8$ کار می‌کند. هر عملوند هشت‌بیتی علامت‌دار و حاصل ضرب
‫شانزده‌بیتی است. بازه‌ی ورودی از ۱۲۸− تا ۱۲۷+، کمینه‌ی حاصل ممکن ۱۶۲۵۶− و بیشینه‌ی
‫حاصل ممکن ۱۶۳۸۴+ است. جدول~\رجوع{جدول:پرونده‌ها} مسئولیت هر پرونده را خلاصه می‌کند.
‫
‫\شروع{لوح}[ht]
‫\تنظیم‌ازوسط
‫\شرح{پیمانه‌ها و مسئولیت آن‌ها در طرح}
‫\شروع{جدول}{|c|p{9.5cm}|}
‫\خط‌پر
‫\سیاه پرونده & \سیاه مسئولیت \\
‫\خط‌پر \خط‌پر
‫\لر{barrel\_ashr.v} & شیفت حسابی راست با فاصله‌ی متغیر \\
‫\لر{skip\_encode.v} & محاسبه‌ی فاصله تا نزدیک‌ترین مرز آینده‌ی بوث \\
‫\لر{booth\_datapath.v} & نگه‌داری ثبات‌ها، جمع یا تفریق، محاسبه‌ی پرش و تولید حاصل \\
‫\لر{booth\_control.v} & ماشین حالت و تولید سیگنال‌های \لر{Load}، \لر{Step} و \لر{Done} \\
‫\لر{booth\_multiplier.v} & اتصال واحد واپایش و مسیر داده در پیمانه‌ی سطح بالا \\
‫\لر{booth\_multiplier\_tb.v} & آزمون‌گر خودکنترل هدفمند و جامع \\
‫\خط‌پر
‫\پایان{جدول}
‫\برچسب{جدول:پرونده‌ها}
‫\پایان{لوح}
‫
‫رابط خارجی شامل \لر{Clk}، بازنشانی فعال‌پایین \لر{RstN}، پالس \لر{Start}، دو عملوند،
‫خروجی \لر{Product} و پالس یک‌چرخه‌ای \لر{Done} است. واحد واپایش دو سیگنال \لر{Load} و
‫\لر{Step} را به مسیر داده می‌فرستد و مسیر داده با \لر{LastStep} پایان محاسبه را اعلام می‌کند.
‫
‫
‫\subsection{اتصال واحد واپایش و مسیر داده}
‫
‫خطوط ۲۲ تا ۴۱ پرونده‌ی \کد{booth\_multiplier.v} اتصال دقیق دو بخش را نشان می‌دهند.
‫شماره‌های کنار فهرست زیر همان شماره‌خط‌های پرونده‌ی اصلی‌اند.
‫
‫\begin{dsdlcodeblock}
‫\begin{verbatim}
‫22      booth_control u_control (
‫23          .Clk(Clk),
‫24          .RstN(RstN),
‫25          .Start(Start),
‫26          .LastStep(LastStep),
‫27          .Load(Load),
‫28          .Step(Step),
‫29          .Done(Done)
‫30      );
‫31
‫32      booth_datapath #(.N(N)) u_datapath (
‫33          .Clk(Clk),
‫34          .RstN(RstN),
‫35          .Load(Load),
‫36          .Step(Step),
‫37          .Multiplicand(Multiplicand),
‫38          .Multiplier(Multiplier),
‫39          .Product(Product),
‫40          .LastStep(LastStep)
‫41      );
‫\end{verbatim}
‫\end{dsdlcodeblock}
‫
‫
‫\section{طراحی مسیر داده}
‫
‫مسیر داده ثبات‌های $A$، $Q$، $Q_{-1}$، $M$، \لر{Remaining} و \لر{Bmap} را نگه
‫می‌دارد. $A$ به جای هشت بیت، نه بیت دارد. این بیت نگهبان برای حالتی ضروری است که در
‫میانه‌ی الگوریتم مقدار $0-(-128)=128$ تولید می‌شود؛ مقدار ۱۲۸+ در قالب علامت‌دار هشت‌بیتی
‫قابل نمایش نیست، اما در نه بیت بدون سرریز نگه‌داری می‌شود.
‫
‫در چرخه‌ی بارگذاری، $A$ و $Q_{-1}$ صفر، $Q$ برابر ضریب، $M$ برابر مضروب و
‫\لر{Remaining} برابر هشت می‌شوند. در هر مرحله‌ی محاسبه، ابتدا عملیات بوث روی $A$ انجام
‫می‌شود؛ سپس ثبات ترکیبی و نقشه‌ی مرزها به اندازه‌ی فاصله‌ی محاسبه‌شده شیفت می‌خورند و همان
‫فاصله از \لر{Remaining} کم می‌شود.
‫
‫
‫\subsection{انتخاب جمع، تفریق و فاصله‌ی شیفت}
‫
‫خطوط ۳۱ تا ۵۵ پرونده‌ی \کد{booth\_datapath.v} منطق تصمیم بوث، ایجاد نقشه‌ی مرز، محاسبه‌ی
‫فاصله و تشخیص آخرین مرحله را پیاده‌سازی می‌کنند.
‫
‫\begin{dsdlcodeblock}
‫\begin{verbatim}
‫31      // Booth pair 01 adds, 10 subtracts, and 00/11 only shift.
‫32      wire DoAdd = (~Q[0]) & Q_1;
‫33      wire DoSub = Q[0] & ~Q_1;
‫34
‫35      // Sign extension preserves negative M values in the wider accumulator.
‫36      wire [N:0] MExt = {M[N-1], M};
‫37
‫38      wire [N:0] AOp = DoSub ? (A - MExt) : DoAdd ? (A + MExt) : A;
‫39
‫40      // A set bit marks a Booth transition: Q[i] XOR Q[i-1], with Q[-1]=0.
‫41      // The map shifts with Q so bit zero always describes the current position.
‫42      wire [N-1:0] BmapInit = Multiplier ^ {Multiplier[N-2:0], 1'b0};
‫43
‫44      // The leading zero lets the encoder return N when no boundary remains.
‫45      wire [DistW-1:0] Raw;
‫46
‫47      skip_encode #(.MapW(N+1)) u_encode (
‫48          .Map({1'b0, Bmap}),
‫49          .Dist(Raw)
‫50      );
‫51
‫52      // Cap the skip at the number of unprocessed multiplier bits.
‫53      wire [DistW-1:0] Dist = (Raw < Remaining) ? Raw : Remaining;
‫54
‫55      assign LastStep = (Dist == Remaining);
‫\end{verbatim}
‫\end{dsdlcodeblock}
‫
‫رابطه‌ی ایجاد نقشه‌ی مرز را می‌توان به صورت زیر نوشت:
‫
‫\شروع{equation}
‫B_i = Q_i \oplus Q_{i-1}, \qquad Q_{-1}=0.
‫\پایان{equation}
‫
‫بیت یک در $B_i$ یعنی بین دو بیت مجاور ضریب تغییر رخ داده و در آن موقعیت یک جمع یا تفریق
‫بوث لازم است. رمزگذار، کوچک‌ترین اندیس یکِ بزرگ‌تر از صفر را انتخاب می‌کند. اندیس صفر به
‫عملیات همین مرحله مربوط است و نباید فاصله‌ی پرش را صفر کند. اگر مرز دیگری وجود نداشته باشد،
‫فاصله‌ی پیش‌فرض برابر $N$ است و سپس با \لر{Remaining} محدود می‌شود.
‫
‫
‫\subsection{شیفت‌دهنده‌ی بشکه‌ای}
‫
‫شیفت‌دهنده‌ی \کد{barrel\_ashr.v} یک مدار ترکیبی پارامتری است. خطوط ۱۲ تا ۱۶ با بررسی
‫بیت‌های مقدار شیفت، جابه‌جایی‌های ۱، ۲، ۴، ۸ و در صورت نیاز ۱۶ بیتی را آبشاری اعمال می‌کنند.
‫
‫\begin{dsdlcodeblock}
‫\begin{verbatim}
‫12      always @* begin
‫13          dout = din;
‫14          for (k = 0; k < AmountW; k = k + 1)
‫15              if (amt[k]) dout = $signed(dout) >>> (1 << k);
‫16      end
‫\end{verbatim}
‫\end{dsdlcodeblock}
‫
‫استفاده از عملگر \کد{>>>} روی مقدار تبدیل‌شده به علامت‌دار، بیت علامت را در شیفت راست حفظ
‫می‌کند. یک نمونه‌ی ۱۸ بیتی برای ${A,Q,Q_{-1}}$ و یک نمونه‌ی ۹ بیتی برای نقشه‌ی مرزها
‫به کار رفته است. ورودی نمونه‌ی دوم با صفر آغاز می‌شود؛ بنابراین شیفت حسابی آن برای نقشه مانند
‫شیفت منطقی عمل می‌کند.
‫
‫
‫\subsection{نمونه‌ی مرحله‌به‌مرحله}
‫
‫برای $M=3$ و $Q=7$، نقشه‌ی اولیه برابر $00001001_2$ است. جدول~\رجوع{جدول:ردیابی}
‫دو مرحله‌ی محاسبه را نشان می‌دهد. این نتیجه با رابطه‌ی $3\times(8-1)=21$ سازگار است.
‫
‫\شروع{لوح}[ht]
‫\تنظیم‌ازوسط
‫\شرح{ردیابی مسیر داده برای ضرب ۳ در ۷}
‫\شروع{جدول}{|c|c|c|c|c|c|}
‫\خط‌پر
‫\سیاه وضعیت & \سیاه زوج بوث & \سیاه عملیات & \سیاه فاصله & \سیاه باقی‌مانده & \سیاه حاصل معتبر \\
‫\خط‌پر \خط‌پر
‫پس از بارگذاری & $10$ & تفریق ۳ & ۳ & ۸ & خیر \\
‫پس از مرحله‌ی ۱ & $01$ & جمع ۳ & ۵ & ۵ & خیر \\
‫پس از مرحله‌ی ۲ & -- & پایان & -- & ۰ & ۲۱ \\
‫\خط‌پر
‫\پایان{جدول}
‫\برچسب{جدول:ردیابی}
‫\پایان{لوح}
‫
‫
‫\section{طراحی واحد واپایش}
‫
‫\subsection{فرض‌های طراحی و انواع حالت‌های ماشین حالت}
‫
‫طرح به صورت پارامتری نوشته شده است، اما مقدار پیش‌فرض پارامتر $N$ برابر ۸ در نظر گرفته
‫می‌شود. همچنین به دلیل استفاده از برش \کد{Multiplier[N-2:0]} در تشکیل نقشه‌ی مرز بوث،
‫فرض شده است که $N\geq 2$ باشد؛ بنابراین پیکربندی یک‌بیتی در محدوده‌ی این پیاده‌سازی نیست.
‫
‫واحد واپایش یک ماشین حالت متناهی چهارحالته با ترتیب
‫\لر{IDLE $\rightarrow$ LOAD $\rightarrow$ CALC $\rightarrow$ REPORT $\rightarrow$ IDLE}
‫و از نوع مور است، زیرا خروجی‌های واپایشی فقط از حالت جاری به دست می‌آیند. چهار نوع حالت آن
‫عبارت‌اند از: حالت انتظار \لر{IDLE}، حالت بارگذاری \لر{LOAD}، حالت محاسبه \لر{CALC} و
‫حالت اعلام نتیجه \لر{REPORT}. در حالت \لر{IDLE} پالس \لر{Start} پذیرفته می‌شود؛ در \لر{LOAD} عملوندها بارگذاری
‫می‌شوند؛ در \لر{CALC} سیگنال \لر{Step} تا فعال شدن \لر{LastStep} یک است؛ و در
‫\لر{REPORT} خروجی \لر{Done} برای یک چرخه فعال می‌شود. خطوط ۱۲ تا ۲۸ پرونده‌ی
‫\کد{booth\_control.v} این رفتار را نشان می‌دهند.
‫
‫\begin{dsdlcodeblock}
‫\begin{verbatim}
‫12      localparam IDLE = 2'd0, LOAD = 2'd1, CALC = 2'd2, REPORT = 2'd3;
‫13
‫14      reg [1:0] State;
‫15
‫16      always @(posedge Clk or negedge RstN) begin
‫17          if (!RstN) State <= IDLE;
‫18          else case (State)
‫19              IDLE: State <= Start ? LOAD : IDLE;
‫20              LOAD: State <= CALC;
‫21              CALC: State <= LastStep ? REPORT : CALC;
‫22              REPORT: State <= IDLE;
‫23          endcase
‫24      end
‫25
‫26      assign Load = (State == LOAD);
‫27      assign Step = (State == CALC);
‫28      assign Done = (State == REPORT);
‫\end{verbatim}
‫\end{dsdlcodeblock}
‫
‫بازنشانی \لر{RstN} ناهمگام و فعال‌پایین است. پالس \لر{Start} هنگام مشغول بودن مدار اثری
‫ندارد و حاصل تا چرخه‌ی بارگذاری بعدی در ثبات‌های مسیر داده باقی می‌ماند.
‫
‫
‫\section{آزمون و ارزیابی}
‫
‫آزمون‌گر \کد{booth\_multiplier\_tb.v} با دوره‌ی ساعت ۱۰ نانوثانیه، ابتدا حالت‌های مرزی و
‫الگوهای هدفمند را اجرا می‌کند و سپس تمام زوج‌های ورودی هشت‌بیتی را می‌آزماید. برای هر زوج،
‫حاصل مدار پس از \لر{Done} با عبارت ضرب علامت‌دار وریلاگ مقایسه می‌شود. خطوط ۱۴۴ تا ۱۶۰
‫پرونده‌ی آزمون‌گر حلقه‌ی جامع و گزارش پایانی را نشان می‌دهند.
‫
‫\begin{dsdlcodeblock}
‫\begin{verbatim}
‫144          for (i = 0; i < 256; i = i + 1)
‫145              for (j = 0; j < 256; j = j + 1)
‫146                  do_mul(i[N-1:0], j[N-1:0]);
‫147
‫148          $display("  %0d products checked against Verilog signed multiply", Runs);
‫149          $display("");
‫150          $display("--- cycle cost per multiply (CALC steps only) ---");
‫151          $display("  plain shift-and-add / 1-bit Booth: %0d steps always", N);
‫152          $display("  this design: min %0d, max %0d, average %0.2f  (%0.1f%% of it)",
‫153                   MinSteps, MaxSteps, SumSteps * 1.0 / Runs,
‫154                   SumSteps * 100.0 / (Runs * N));
‫155
‫156          $display("");
‫157          if (Errors == 0)
‫158              $display("========== PASS: booth_multiplier clean ==========");
‫159          else
‫160              $display("========== FAIL: %0d error(s) ==========", Errors);
‫\end{verbatim}
‫\end{dsdlcodeblock}
‫
‫
‫\subsection{نتایج شبیه‌سازی}
‫
‫منابع فعلی با \لر{Icarus Verilog} و گزینه‌های \کد{-Wall} و
‫\کد{-g2005-sv} کامپایل و اجرا شدند. هر $256\times256=65{,}536$ زوج ورودی با
‫مرجع مقایسه شد و پیام نهایی \کد{PASS} بود. جدول~\رجوع{جدول:نتایج هدفمند} بخشی از
‫حالت‌های هدفمند و تعداد مراحل \لر{CALC} آن‌ها را نشان می‌دهد.
‫
‫\شروع{لوح}[ht]
‫\تنظیم‌ازوسط
‫\شرح{نمونه‌هایی از نتایج آزمون هدفمند}
‫\شروع{جدول}{|c|c|c|}
‫\خط‌پر
‫\سیاه ضرب & \سیاه حاصل & \سیاه مراحل محاسباتی \\
‫\خط‌پر \خط‌پر
‫$0\times0$ & ۰ & ۱ \\
‫$3\times7$ & ۲۱ & ۲ \\
‫$-3\times5$ & ۱۵− & ۴ \\
‫$-128\times-128$ & ۱۶۳۸۴ & ۲ \\
‫$-128\times127$ & ۱۶۲۵۶− & ۲ \\
‫$85\times-86$ & ۷۳۱۰− & ۸ \\
‫$5\times-1$ & ۵− & ۱ \\
‫\خط‌پر
‫\پایان{جدول}
‫\برچسب{جدول:نتایج هدفمند}
‫\پایان{لوح}
‫
‫کمینه‌ی مراحل محاسباتی ۱، بیشینه ۸ بود. روش شیفت و جمع یا بوث یک‌بیتی
‫برای ورودی هشت‌بیتی همواره ۸ مرحله نیاز دارد. بنابراین نسبت میانگین مراحل طرح حاضر برابر است با:
‫
‫\شروع{equation}
‫\frac{4.50}{8}\times100 = 56.25\%.
‫\پایان{equation}
‫
‫این مقایسه فقط
‫چرخه‌های \لر{CALC} را در نظر می‌گیرد؛ چرخه‌های ثابت بارگذاری و گزارش در هر دو نتیجه‌ی این
‫محاسبه وارد نشده‌اند.
‫
‫\شروع{شکل}[H]
‫\تنظیم‌ازوسط
‫\centerimg{terminal1.png}{16cm}
‫
‫\شرح{خروجی تست}
‫\برچسب{شکل:نمای RTL}
‫\پایان{شکل}
‫
‫
‫\شروع{شکل}[ht]
‫\تنظیم‌ازوسط
‫\centerimg{main-wave-1.png}{16cm}%
‫\شرح{شکل موج یک اجرای کامل ضرب‌کننده در شبیه‌سازی}
‫\برچسب{شکل:شکل موج}
‫\پایان{شکل}
‫
‫\شروع{شکل}[ht]
‫\تنظیم‌ازوسط
‫\centerimg{main-wave-2.png}{16cm}%
‫\شرح{ادامه شکل موج یک اجرای کامل ضرب‌کننده در شبیه‌سازی}
‫\برچسب{شکل:شکل موج}
‫\پایان{شکل}
‫
‫
‫\section{سنتز \lr{Quartus}}
‫
‫\شروع{شکل}[H]
‫\تنظیم‌ازوسط
\centerimg{quartus1.jpg}{16cm}

‫\شرح{ گزارش کامپایل موفق طرح با مصرف ۱۶۰ عنصر منطقی، ۴۱ ثبات و ۳۶ پایه.}
‫\برچسب{شکل:نمای RTL}
‫\پایان{شکل}
‫
‫\شروع{شکل}[H]
‫\تنظیم‌ازوسط
‫\centerimg{quartus2.jpg}{16cm}
‫
‫\شرح{ نمای سطح بالای اتصال مسیر داده و واحد کنترل ضرب‌کننده‌ی بوث.}
‫\برچسب{شکل:نمای RTL}
‫\پایان{شکل}
‫
‫\شروع{شکل}[H]
‫\تنظیم‌ازوسط
‫\centerimg{quartus3.jpg}{16cm}
‫
‫\شرح{ نمای فناوری مدار سنتزشده شامل منطق داخلی و بافرهای ورودی و خروجی.}
‫\برچسب{شکل:نمای RTL}
‫\پایان{شکل}
‫
‫
‫\section{نتیجه‌گیری}
‫
‫در این آزمایش یک ضرب‌کننده‌ی علامت‌دار بوث با مسیر داده و واحد واپایش مستقل طراحی شد. نقشه‌ی
‫مرزها و شیفت‌دهنده‌ی بشکه‌ای امکان مصرف چند بیت ضریب در یک مرحله را فراهم کردند؛ در نتیجه
‫میانگین مراحل محاسباتی برای تمام زوج‌های هشت‌بیتی از ۸ به ۵۰٫۴ کاهش یافت. بیت نگهبان انباره
‫نیز صحت حالت‌های مرزی مکمل دو را حفظ کرد. آزمون جامع همه‌ی زوج ورودی بدون خطا
‫پایان یافت و صحت تابعی مدل وریلاگ را تأیید کرد.
‫
‫
‫\پایان{شمارش}
‫